Read in Gauge Theory Gravity: the metric $g$ below is the derived metric of the gravitational field $(\hb, \Omega)$ on flat spacetime, the manifold is the invertible region, and the tractor connection is the flat-space $\Cl(2,4)$ gauge field of Theorem \claimref{synthesis.gtg-tractor-gauge-field-reproduces-tractor-connection}; the index computation below is its reading in the splitting of $g$ (Fact \claimref{synthesis.imported-gtg-dictionary}).

\emph{Setting.} Tractors in the splitting of $g$ are $(\sigma, \mu_b, \rho)$ with the form $\ip{(\sigma,\mu,\rho)}{(\sigma',\mu',\rho')} = \sigma\rho' + \rho\sigma' + g^{bc}\mu_b\mu'_c$, and by Fact \claimref{interior.imported-tractor-connection} the connection along $v$ is $\nabla_v(\sigma,\mu,\rho) = (\nabla_v\sigma - \mu(v),\ \nabla_v\mu + v^\flat\rho - \Sch(v,\cdot)\sigma,\ \nabla_v\rho + \Sch(v,\mu^\sharp))$, with $\Sch = \tfrac12(\Ric - \tfrac{R}{6}g)$ and $J = R/6$ (Lemma \claimref{interior.metric-parallelism-trace-schouten-condition}(ii)). For $I_g = (1,0,0)$ and $X = (0,0,1)$ this gives $\nabla_v I_g = (0, -\Sch(v,\cdot), 0)$ and $\nabla_v X = (0, v^\flat, 0)$, so the classes against $\mu$-slot vectors are $-\Sch(v,w)$ and $g(v,w)$, and by Theorem \claimref{synthesis.einstein-condition-is-infinity-vacuum-incidence}(2) they depend only on the vacua at $x$.

\emph{Why only the canonical classes.} For a field $W$ of maximal isotropic subspaces and sections $w, w'$ of it, the class $\ip{\nabla_v w}{w'}$ is tensorial in $w$ and $w'$ (incidence theorem, proof of (2)), but for $w$ not a multiple of $I_g$ or $X$ it depends on the field's first-order variation at $x$: two fields that agree at $x$ and differ to first order (for example by a rotation of the plane $A_W$ along $v$ by an element of the Lie algebra of the tractor orthogonal group) give classes that differ by the pairing of that element. That variation is a choice of extension, not a function of $g$, so a tensor meant to be determined by the geometry at $x$ may depend only on the classes of the canonical sections; this is the hypothesis of the statement, and it is argued here, not derived from a premise.

\emph{(1) Form.} The stabiliser of the two null vectors $I_g$, $X$ with $\ip{I_g}{X} = 1$ fixes the hyperbolic plane they span and acts on its orthogonal complement, the $\mu$ slot with the form $g$ of signature $(1,3)$, as $O(1,3)$; on the classes it acts by $\Sch \mapsto \Lambda^{T}\Sch\Lambda$. Let $K$ be the space of linear maps of symmetric two-tensors commuting with the action of the five elements $B$ (boost in the $(0,1)$ plane with $\cosh = 5/4$, $\sinh = 3/4$), $R_{12}$ (rotation with $\cos = 3/5$, $\sin = 4/5$), $R_{23}$ (rotation with $\cos = 5/13$, $\sin = 12/13$), and the reflections of $e_0$ and of $e_1$; each preserves $\eta = \mathrm{diag}(1,-1,-1,-1)$. The exact computation in the check gives $\dim K = 2$, and $K$ contains the identity and $\Sch \mapsto (\tr_g\Sch)\,g$, which commute with all of $O(1,3)$. The commutant of $O(1,3)$ is contained in $K$ (it commutes in particular with the five elements) and contains these two independent maps, so it equals their span. Hence $L(\Sch) = a\Sch + bJg$ and $\Theta = a\Sch + bJg + cg$.

\emph{(2) Dichotomy.} For real null $\ell$, $g(\ell,\ell) = 0$ gives $\Theta(\ell,\ell) = a\Sch(\ell,\ell)$, and $\omega_\ell(I_g,\ell) = -\Sch(\ell,\ell)$, so the clause is $(1 + \kappa a)\Sch(\ell,\ell) = 0$. If $a = -1/\kappa$ every $\Sch$ satisfies it. Otherwise it says $\Sch(\ell,\ell) = 0$ for every real null $\ell$, which by the incidence theorem (4) (the symmetric forms vanishing on the null cone are the multiples of $g$; in the check, rank nine on eleven rational null vectors with kernel spanned by $\eta$) is $\Sch = \tfrac{J}{4}g$. Then $\Theta = (\tfrac{a}{4} + b)Jg + cg$, pure trace, and $\Theta(\ell,\ell) = 0$.

\emph{(3) Conservation.} The contracted Bianchi identity $\nabla^a\Ric_{ab} = \tfrac12\nabla_b R$ holds in either overall sign convention for $\Ric$ (both sides change sign together), so $\nabla^a\Sch_{ab} = \tfrac12(\tfrac12 - \tfrac16)\nabla_b R = \tfrac16\nabla_b R = \nabla_b J$. With $a, b, c$ constants, $\nabla^a\Theta_{ab} = a\nabla_b J + b\nabla_b J$. There are metrics with $\nabla J \neq 0$ at a point (the check exhibits one, with first derivatives of $g$ vanishing at the origin and generic second and third derivatives, and verifies $\nabla^a\Sch_{ab} = \nabla_b J$ there exactly: since $\Gamma(0) = 0$, the terms of the Riemann tensor quadratic in $\Gamma$ have vanishing value and first derivative at the origin, and $\nabla^a\Sch_{ab}(0) = \eta^{ac}\partial_c\Sch_{ab}(0)$). So conservation for every metric is $a + b = 0$ (the linear case of Lovelock's theorem, which says a conserved symmetric tensor built linearly from the curvature in four dimensions is a multiple of the Einstein tensor plus a multiple of $g$; cited from memory, not used beyond this computation), in case (a) $b = 1/\kappa$, and
\[
\Theta = -\tfrac{1}{\kappa}(\Sch - Jg) + cg, \qquad \Sch - Jg = \tfrac12\Ric - \tfrac{R}{12}g - \tfrac{R}{6}g = \tfrac12\bigl(\Ric - \tfrac12 R\,g\bigr),
\]
so $\Theta = -\tfrac{1}{8\pi G}(\Ric - \tfrac12 Rg) + cg$ and $\Ric + 8\pi G\,\Theta = (\tfrac12 R + 8\pi G c)g$ for every metric. In the store's convention $\Ric = -\Ric^{\circ}$ and $R = -R^{\circ}$, with $\Ric^{\circ}$ the sphere-positive Ricci tensor, so $\Theta = \tfrac{1}{8\pi G}(\Ric^{\circ} - \tfrac12 R^{\circ}g) + cg$: the Einstein tensor over $8\pi G$ plus a cosmological term. Clause (7) of the incidence theorem with $T := \Theta$ then holds with $h = \tfrac12 R + 8\pi G c$ for every metric; it constrains nothing.

\emph{(4) Consequence.} By (1) a first-order vacuum-field tensor is determined at $x$ by $\Sch(x)$ and three numbers. In case (a) with $a, b, c$ fixed, $\Theta(x)$ is a function of $\Sch(x)$, so it cannot take two values with different $\Theta(\ell,\ell)$ at one $\Sch$; in case (b) the clause is satisfiable only with $\Theta(\ell,\ell) = 0$. A source with either property in the statement is therefore not of this form, and clause (7) with such a source relates two independent data only through an added hypothesis. The argument concerns first-order classes only; it does not exclude a source that is not a function of them, for example an independent $T(k,k)$ given as the density of the generator of null translations of a field theory on the patch, which is a datum of the kind (4) asks for.

\emph{Remark (higher orders; argued, not used).} Second-order data of the vacuum field are the tractor curvature, which carries the Weyl tensor and the Cotton tensor, and $\nabla\Sch$. Since $-1 \in O(1,3)$ acts on tensors of rank $k$ by $(-1)^{k}$, no nonzero equivariant linear map takes the rank-three tensors $\nabla\Sch$ and the Cotton tensor to symmetric two-tensors, and by Schur's lemma none takes the Weyl tensor (an irreducible, trace-free representation inequivalent to the two summands of $\mathrm{Sym}^2$) to them. At the next order $\nabla\nabla\Sch$ gives $\square\Sch$, $\nabla\nabla J$ and $g\,\square J$, and a source built from them makes clause (7) a higher-derivative equation for the metric or an identity, not Einstein's equation with an independent source. The quadratic option, the probability of an excitation, is quadratic in the curvature (a Bel-Robinson type quantity, as argued in earlier work on this question), so it is not linear in a matter source either.

\emph{The check.} With exact rational arithmetic: the classes of $I_g$ and $X$ along three directions against the four frame vectors for a generic rational $\Sch$, and that the $X$ classes do not depend on $\Sch$; the commutant computation of (1) (a $100$-unknown linear system, kernel of dimension two, containing the identity and the trace map); for three triples $(a,b,c)$ with $1 + \kappa a \neq 0$ (with $\kappa = 3$ standing for $4\pi G$; every step is linear in $\kappa$) the rank-nine own-ray system with kernel spanned by $\eta$, and that $\Theta$ is then pure trace with $\Theta(\ell,\ell) = 0$ on eleven null vectors, and for $a = -1/\kappa$ that the system vanishes identically; and on a metric with vanishing first derivatives and generic second and third derivatives at the origin, the contracted Bianchi identity $\nabla^a\Sch_{ab} = \nabla_b J$ with $\nabla J \neq 0$, the divergence formula $(a+b)\nabla J$ for two pairs $(a,b)$, and the identity $\Ric + 2\kappa\Theta = (\tfrac12 R + 2\kappa c)g$ for the conserved tautological branch.
